\section{Materials and methods}
\subsection{Study area, inventory and evidence sources}
Taxila is a serial World Heritage property in Punjab, Pakistan, inscribed under reference 139 \citep{unesco_taxila}. The inventory contained 18 component records. Seventeen had usable point coordinates; Saraikala (139-002) remained in the documentary inventory but was excluded from coordinate-dependent analysis because its location was unresolved. Spectral neighbourhoods used WGS 84/UTM zone 43N (EPSG:32643); terrain extraction used local metric approximations on the geographic elevation grid. Point-centred circles of 250, 500 and 1,000 m defined analytical neighbourhoods, with 500 m used for primary reporting. They are neither official property polygons nor statutory or UNESCO buffer zones (Fig.~\ref{fig:study}).

\begin{figure}[tbp]\centering
\includegraphics[width=\textwidth]{figure_01_study_area_context_map.pdf}
\caption{Taxila study-area context and the official 18-component inventory. Seventeen components are mapped; Saraikala is retained in the inventory without an inferred coordinate. Dark-red circles indicate 500 m analytical neighbourhoods. Esri World Imagery provides contextual imagery only; it was not used to calculate spectral, climate or terrain indicators. Boundaries: Natural Earth and geoBoundaries; road and museum: OpenStreetMap contributors (ODbL). Map lines delineate study areas and do not necessarily depict accepted national boundaries. }\label{fig:study}
\end{figure}

Landsat Collection 2 Level-2 surface reflectance supplied five post-monsoon three-year composites: 2003--2005, 2008--2010, 2013--2015, 2018--2020 and 2023--2025, abbreviated E2004 to E2024. They contained 11, 10, 8, 8 and 18 retained scenes, respectively. Landsat 7 supplied the first two epochs, Landsat 8 the middle two, and Landsat 8/9 the final epoch. Source records preserve acquisition dates, sensor identities and quality information. Cross-sensor and masking limitations were considered explicitly \citep{usgs_landsat_c2,roy2016landsat}.

Daily Open-Meteo ERA5-Seamless data covered 1991--2025 \citep{openmeteo_archive}. Component queries resolved to four unique weather cells, each containing 12,784 daily records. NASA POWER supplied a comparison product over the same period \citep{nasa_power}. Neither was treated as station truth. Terrain came from the one-arc-second N33E072 elevation tile distributed through AWS Terrain Tiles. ESA WorldCover 2020/2021 consensus \citep{esa_worldcover} supplied six proxy land-cover classes for an E2019 feature table. Table~\ref{tab:data} summarises the role of each source. Figure~\ref{fig:framework} separates the paths through which these inputs inform interpretation.

\begin{table}[tbp]\centering\small
\caption{Evidence sources and permitted analytical roles.}\label{tab:data}
\begin{tabular}{p{.22\linewidth}p{.30\linewidth}p{.38\linewidth}}\toprule
Source & Support & Analytical role \\\midrule
Landsat C2 L2 & 55 scenes; five epochs; 30 m grid & Relative spectral pressure and endpoint convergence \\
Elevation tile & Approximately 30 m & Terrain and drainage susceptibility \\
Open-Meteo & Daily, 1991--2025; four cells & Regional trends and acquisition-window context \\
NASA POWER & Daily, 1991--2025 & Weather-product sensitivity \\
WorldCover consensus & 2020/2021 labels matched to E2019 features & Weakly supervised land-cover benchmark \\
UNESCO component inventory & 18 records; 17 mapped points & Component identifiers and neighbourhood centres \\\bottomrule
\end{tabular}
\end{table}

\begin{figure}[p]\centering
\includegraphics[width=\textwidth]{figure_02_chip_framework.pdf}
\caption{CHIP framework linking observation support, mathematical transformations, experiments and conservation interpretation. Solid arrows carry analytical inputs; dashed arrows carry interpretation. Climate retains native support and land-cover evaluation remains a separate diagnostic. Fixed and reconstructed spectral percentile populations are distinguished. Counts describe executed designs. Codex assisted this editable vector schematic.}\label{fig:framework}
\end{figure}


\subsection{Spectral transformations and common-support convergence}
Cloud, shadow and invalid observations were masked before compositing on a fixed 30 m grid. With scaled reflectance $\rho_b$ in band $b$, the indices were
\begin{equation}\label{eq:indices}
\mathrm{NDVI}=\frac{\rho_{\mathrm{NIR}}-\rho_{\mathrm{red}}}{\rho_{\mathrm{NIR}}+\rho_{\mathrm{red}}},\quad
\mathrm{NDBI}=\frac{\rho_{\mathrm{SWIR1}}-\rho_{\mathrm{NIR}}}{\rho_{\mathrm{SWIR1}}+\rho_{\mathrm{NIR}}},\quad
\mathrm{MNDWI}=\frac{\rho_{\mathrm{green}}-\rho_{\mathrm{SWIR1}}}{\rho_{\mathrm{green}}+\rho_{\mathrm{SWIR1}}}.
\end{equation}
These ratios follow their original definitions \citep{tucker1979ndvi,zha2003ndbi,xu2006mndwi}. NDBI responds to bare soil as well as built surfaces; its adverse direction is a screening assumption, not an attribution of urban expansion.

For component $i$, epoch $e$, radius $r$ and factor $k$, let $x_{ierk}$ denote the neighbourhood median. The oriented empirical percentile is
\begin{equation}\label{eq:percentile}
z_{ierk}=\frac{\operatorname{rank}_{\mathcal I_r}(s_kx_{ierk})-\tfrac12}{N_r},
\qquad s_N=s_W=-1,\quad s_B=+1.
\end{equation}
Ranks ascend with average ties; $N_r=|\mathcal I_r|=85$ component--epoch observations at each radius ($17\times5$). Terrain percentiles use the 17 static component summaries. Spatial reconstruction experiments instead rerank the 17 simultaneously reconstructed E2024 summaries. Reconstruction differences therefore include a reference-population effect. Percentiles describe relative evidence within their stated population, not transferable physical thresholds.

Endpoint differences $\Delta I_k(p)=I_{k,\mathrm{E2024}}(p)-I_{k,\mathrm{E2004}}(p)$ used only the common valid-cell set $\Omega$. With $Q_k(q)$ denoting its empirical $q$-quantile,
\begin{equation}\label{eq:convergence}
\begin{aligned}
C(p)&=\mathbf1[\Delta I_N(p)\leq Q_N(.2)]+\mathbf1[\Delta I_W(p)\leq Q_W(.2)]
       +\mathbf1[\Delta I_B(p)\geq Q_B(.8)],\\
\pi_{\geq h}&=|\Omega|^{-1}\sum_{p\in\Omega}\mathbf1[C(p)\geq h],\qquad h\in\{2,3\}.
\end{aligned}
\end{equation}
The indicator function $\mathbf1[\cdot]$ equals one when its condition holds. These overlapping, data-derived tails quantify descriptive convergence; they are neither independent tests nor measurements of damaged area.

\subsection{Terrain representation and hierarchical decision rule}
D8 routing directs each elevation cell to its steepest strictly lower neighbour; sinks remain terminal. Slope $\theta$ derives from metric elevation gradients. The contributing-cell count $A_c$ and geometric-mean cell width $h$ define the numerical wetness proxy
\begin{equation}\label{eq:wetness}
H(p)=\log\!\left[\frac{(A_c(p)+1)\,h/(1\,\mathrm{m})}
 {\tan(\max\{\theta(p),0.05^\circ\})+0.01}\right].
\end{equation}
This regularised topographic proxy follows the contributing-area/slope rationale of hydrological terrain analysis \citep{beven1979topmodel,moore1991terrain}. Drainage comprises cells at or above the 99th percentile of accumulation. Component factors are the 90th percentiles of slope and $H$, and median Euclidean drainage distance. Higher slope and wetness and shorter distance are oriented towards greater susceptibility. Routing approximates topographic convergence, not discharge, and omits unresolved walls, culverts and local drains.

With $S,H,D$ indexing slope, wetness and drainage proximity, landscape pressure $L$, terrain susceptibility $T$ and priority $P$ are
\begin{equation}\label{eq:priority}
\begin{aligned}
L_{ier}&=\tfrac14z_{ierN}+\tfrac14z_{ierB}+\tfrac12z_{ierW},\\
T_{ir}&=\tfrac13(z_{irS}+z_{irH}+z_{irD}),\\
P_{ier}(w)&=wL_{ier}+(1-w)T_{ir},\qquad P_{ier}^{\mathrm{ref}}=P_{ier}(0.5).
\end{aligned}
\end{equation}
Grouping NDVI and NDBI limits their combined weight relative to MNDWI without claiming statistical independence. Equal domain weights are declared choices, not fitted conservation parameters. This hierarchical construction connects GIS multicriteria assessment \citep{malczewski2006gis} to explicitly testable sensitivity assumptions \citep{chen2010sensitivity,feizizadeh2014uncertainty}. A regional climate term remains contextual: adding the same value to every component through an equal-weight integrated score is an affine transformation of $P^{\mathrm{ref}}$ and cannot alter the within-epoch ordering.

\subsection{Structural experiments and exact weight transitions}
The hierarchy was compared with an equal six-factor mean, landscape-only and terrain-only baselines. Factor ablations removed one term and renormalised retained weights. Spearman correlation, Kendall concordance, top-five overlap and maximum rank displacement describe ordering changes.

Three radii, three landscape weights ($0.25,0.50,0.75$) and three terrain formulations produce 27 scenarios. Terrain alternatives comprise equal factors, inverse-correlation weights and the reduced slope--drainage mean. For the correlation-adjusted case, $a_k\propto[\sum_\ell|\rho_{k\ell}|]^{-1}$ with $\sum_ka_k=1$, where $\rho$ is the Spearman matrix of the three terrain factors, including its diagonal. Continuous-weight experiments sample $w\sim\mathrm{Beta}(\alpha,\alpha)$ for $\alpha\in\{0.5,1,2\}$, with 50,000 draws each. This evaluates stated preference distributions rather than inferring stakeholders' preferences, consistent with rank-acceptability analysis \citep{lahdelma2001smaa}.

An additional deterministic analysis resolves every pairwise score crossing. For E2024 components $i,j$ at fixed radius and equal-factor terrain, solving $P_i(w)=P_j(w)$ gives
\begin{equation}\label{eq:crossing}
w^*_{ij}=\frac{T_j-T_i}{(L_i-T_i)-(L_j-T_j)}.
\end{equation}
Only nonzero denominators and crossings in $(0,1)$ are retained. Sorting these crossings partitions the weight interval into regions of constant rank order; midpoint evaluation identifies the leader in each region. Boundary ties retain component-inventory order. Figure~\ref{fig:scales} combines fixed ranks, a three-dimensional rank response and exact leadership intervals. Radii remain discrete; no interpolated spatial-resolution surface is implied.

\subsection{Spatial, positional and joint uncertainty propagation}
At each block size (90, 150, 300 m), 2,000 spatial draws used shared exponential weights. If $m_{ibk}$ is the median of factor $k$ in block $b$ intersecting component $i$, and $n_{ib}$ its valid-cell count,
\begin{equation}\label{eq:block}
E_b^{(d)}\sim\operatorname{Exp}(1),\quad v_{ib}^{(d)}=n_{ib}E_b^{(d)},\quad
\widetilde x_{ik}^{(d)}=\operatorname{wmed}_{b}\{m_{ibk};v_{ib}^{(d)}\}.
\end{equation}
The weighted median is the first ordered block median reaching half the cumulative weight. Sharing $E_b^{(d)}$ across overlapping neighbourhoods preserves their resampling dependence; cell counts retain unequal spatial support. Terrain remains fixed in this experiment.

For positional sensitivity, independent uniform-area displacements are
\begin{equation}\label{eq:displacement}
\boldsymbol\delta_i=d_{\max}\sqrt{U_i}(\cos\phi_i,\sin\phi_i),\quad
U_i\sim U(0,1),\quad \phi_i\sim U(0,2\pi).
\end{equation}
Spectral and terrain summaries were re-extracted for 2,000 draws at each $d_{\max}\in\{30,60,120\}$ m. An undisplaced re-extraction supplies the appropriate control. Equation~\eqref{eq:percentile} distinguishes this endpoint reference from the fixed summary-table reference.

Joint simulation generates 500 spatial states, each with sampled block size, shared block weights and independent displacement within 120 m. Ten decision settings per state sample $w\sim U(0,1)$ and one terrain formulation. For resulting ordinal ranks $R_i^{sd}$, top-$k$ frequency is
\begin{equation}\label{eq:acceptability}
\widehat b_i(k)=\frac{1}{500\times10}\sum_{s=1}^{500}\sum_{d=1}^{10}
 \mathbf1[R_i^{sd}\leq k].
\end{equation}
Median ranks and central 95\% intervals summarise these 5,000 nested score vectors. They are conditional sensitivity distributions, not 5,000 independent maps or posterior probabilities of deterioration. The same frequency definition over 27 structural scenarios defines a core tier at $b_i(3)\geq0.5$ and a conditional tier outside it at $b_i(5)\geq0.5$.

\subsection{Climate trends and acquisition-matched forcing}
The regional daily series averages four unique weather cells equally before annual aggregation. For annual values $Y_t$, the Theil--Sen estimator is
\begin{equation}\label{eq:sen}
\widehat\beta=\operatorname{median}_{u>t}\frac{Y_u-Y_t}{u-t}.
\end{equation}
After robust trend fitting, five-year residual blocks were resampled and the fitted trend restored to obtain 2,000 slope replicates \citep{sen1968slope,yue2002autocorrelation}. Null replicates add resampled residuals to a constant baseline. The two-sided probability is $(1+\sum_b\mathbf1[|\beta_b^0|\geq|\widehat\beta|])/(2,001)$; Benjamini--Hochberg adjustment covers 12 annual metrics \citep{benjamini1995fdr}. Inference assumes adequate representation of short-range residual dependence and remains conditional on the weather product.

For acquisition dates $a\in\mathcal A_e$, daily precipitation $W(t)$ and antecedent window $\ell$,
\begin{equation}\label{eq:lags}
\overline W_e(\ell)=\frac{1}{|\mathcal A_e|}\sum_{a\in\mathcal A_e}
 \sum_{u=1}^{\ell}W(a-u),\qquad \ell\in\{7,30,90,180,365\}\ \mathrm{days}.
\end{equation}
Temperature uses window means. Product comparisons match dates and periods. Five composite epochs provide only five temporal replicates for exploratory weather--spectral associations; no weather field was downscaled to 30 m.

\subsection{Buffered model selection and probability diagnostics}
The weakly supervised benchmark contained 20,749 class-stratified samples, six reflectance bands and four indices. Its outer stripe contained 4,188 samples in 20 blocks; 4,467 samples in adjacent block columns were excluded, leaving 12,094 development samples in 60 blocks. Three contiguous inner validation bands each excluded an adjacent block row from training. This design targets geographical transfer \citep{roberts2017cv,meyer2018target}, with class-stratified agreement interpreted separately from population-area map accuracy \citep{olofsson2014accuracy,stehman2019accuracy}.

Seven families---logistic regression, random forest, extra trees, histogram gradient boosting, MLP, XGBoost and CatBoost---each received three predeclared configurations \citep{breiman2001rf,pedregosa2011sklearn,chen2016xgboost,prokhorenkova2018catboost}. The resulting 63 inner fits selected family $m$ and configuration $\eta$ through
\begin{equation}\label{eq:selection}
\begin{aligned}
(\widehat m,\widehat\eta)&=\arg\max_{m,\eta}\frac13\sum_{v=1}^{3}
 F_{1,\mathrm{macro}}(y_v,\widehat f_{m,\eta}^{(-v)}(X_v)),\\
F_{1,\mathrm{macro}}&=\frac16\sum_{c=1}^{6}
 \frac{2\mathrm{TP}_c}{2\mathrm{TP}_c+\mathrm{FP}_c+\mathrm{FN}_c}.
\end{aligned}
\end{equation}
Undefined class terms receive zero. Water was absent from one inner training fold; the six-class denominator was retained and unseen-class probability set to zero where necessary. Standardisation was fitted within training folds and early stopping disabled. All families were refitted on development data before outer prediction; selection used only inner scores.

Alongside macro-F1, fixed-six-class mean recall and overall agreement, probability quality used
\begin{equation}\label{eq:probability}
\mathcal L=-\frac1n\sum_{j=1}^n\log p_{j,y_j},\qquad
\mathrm{BS}=\frac1n\sum_{j=1}^n\sum_{c=1}^{6}(p_{jc}-\mathbf1[y_j=c])^2.
\end{equation}
Log loss and the unnormalised multiclass Brier score assess probabilistic predictions \citep{gneiting2007scoring}. Ten equal-width confidence bins define top-label calibration error as $\sum_b(n_b/n)|\mathrm{accuracy}_b-\mathrm{confidence}_b|$. A scalar $\tau\in[0.25,4]$ minimised development out-of-fold log loss, using
\begin{equation}\label{eq:temperature}
p_{jc}^{(\tau)}=\frac{\exp[\log(\max\{p_{jc},10^{-12}\})/\tau]}
 {\sum_{h=1}^{6}\exp[\log(\max\{p_{jh},10^{-12}\})/\tau]}.
\end{equation}
The transform follows temperature scaling \citep{guo2017calibration}; its transfer was tested on outer probabilities. Two thousand shared test-block resamples estimated paired macro-F1 differences. Bands-only, indices-only, location-only, shuffled-label and class-prior controls tested information content; three seeds assessed refit sensitivity.

The outer stripe had been examined previously, making this a retrospective comparison. WorldCover consensus labels also differ temporally from E2019 predictors. These experiments therefore evaluate land-cover proxy agreement without establishing independent monument-condition accuracy.

\subsection{Computational implementation}
Model experiments ran in Google Colab and independently on Windows with recorded package versions, spatial splits, configurations and per-sample probabilities. Derived inputs and deterministic figure/table builders support reproduction \citep{wilson2017practices,nust2018reproducible}. Scene manifests document acquisition and preprocessing. The exact weight analysis uses recorded component summaries without refitting their percentile transforms.
